\documentclass[11pt]{article}
\usepackage[final]{acl}
\usepackage{times}
\usepackage{latexsym}
\usepackage{booktabs}
\usepackage{multirow}
\usepackage[T1]{fontenc}
\usepackage[utf8]{inputenc}
\usepackage{microtype}
\usepackage{inconsolata}
\usepackage{xeCJK}
\setCJKmainfont{Songti SC}

\title{Which Chinese Standard? Mainland and Taiwan Grade 65\% of\\
       Shared Vocabulary Differently, and No Alignment Fixes It}

\author{Alvaro Serrano \\
  \textit{Independent researcher} \\
  \texttt{alvaro.serrano.gp@gmail.com} \\
  \texttt{orcid.org/0009-0006-4701-9026}}

\begin{document}
\maketitle

\begin{abstract}
Chinese has two national proficiency ecosystems. The mainland publishes a
grading standard and an examination syllabus; Taiwan publishes TBCL and TOCFL.
Learners, publishers and researchers move between them constantly, but the
published comparisons are exam-choice guidance --- level equivalency tables and
headline vocabulary counts --- and no word-by-word comparison of the
inventories appears to exist. We provide one. Of the 8{,}318 words TBCL and the
2025 mainland syllabus share, \textbf{64.7\% sit at different tiers}; across the
character inventories the figure is 55.2\%. Both far exceed the 41.5\%
disagreement between the mainland's \emph{own} two documents
\citep{serrano2026whichhsk}, which our pipeline reproduces exactly as a control.
We also report a claim we cannot make. The obvious next question --- which side
grades harder --- appears to have an emphatic answer, 4{,}783 words against 598.
It does not survive: TBCL levels 1 and 2 sit below CEFR A1 while the mainland
scale has no sub-A1 tier, and shifting TBCL by a single level reverses the
direction to 1{,}234 against 2{,}421. The disagreement rate holds under every
alignment tested, between 43.9\% and 96.2\%; the direction is an artifact of a
correspondence no issuing body has published. We release the extraction and
comparison code, and a validated machine-readable extraction procedure for
TBCL, for which none appears to have existed.
\end{abstract}

\section{Two ecosystems, never compared}

A learner who studies for HSK and then sits TOCFL, a publisher levelling a
series for both markets, or a researcher assembling a Chinese resource, all face
the same question: do the two standards mean the same thing by a level? The
literature does not answer it. What exists is comparison for the purpose of
choosing an examination --- equivalency tables asserting that HSK 4 is
``about'' TOCFL 3, and totals for how many words each requires. Those are
statements about scales, not about vocabulary.

The question matters more than it did. The mainland's 2025 examination syllabus
took effect in July 2026, and its arrival prompted a widely repeated claim that
mainland word counts have now overtaken Taiwan's. That claim turns out to be
true and false at once, depending on a distinction nobody making it draws
(\S\ref{sec:counts}).

We compare the inventories directly, in the manner of
\citet{serrano2026whichhsk}, which measured the same disagreement between the
mainland's two documents.

\section{The four documents}
\label{sec:counts}

\begin{table}[t]
\centering\small
\begin{tabular}{@{}llrr@{}}
\toprule
& \textbf{Document} & \textbf{Words} & \textbf{Chars} \\
\midrule
\multirow{2}{*}{PRC}
 & GF0025-2021 standard  & 10{,}916 & 3{,}000 \\
 & 2025 exam syllabus    & 10{,}896 & 3{,}088 \\
\addlinespace
\multirow{2}{*}{TW}
 & TBCL standard         & 14{,}452 & 3{,}100 \\
 & TOCFL 8000 exam list  &  7{,}517 & --- \\
\bottomrule
\end{tabular}
\caption{The four inventories. Taiwan figures are as published; TBCL's word
list carried 14{,}425 entries before the April 2025 revision.}
\label{tab:docs}
\end{table}

Table~\ref{tab:docs} settles the overtaking claim, and shows why it was
ambiguous. Compared exam to exam, the mainland syllabus is 45\% larger than
TOCFL's list and the claim holds. Compared standard to standard, TBCL is 32\%
larger than the mainland standard and the claim fails. The two comparisons give
opposite answers, and the sources making the claim do not say which pair they
mean --- the same failure of document identity that
\citet{serrano2026whichhsk} documents within the mainland, now on the
cross-strait axis.

\paragraph{Extracting TBCL.} No machine-readable extraction of TBCL appears to
exist. NAER publishes the lists as spreadsheets and asserts rights over them, so
we release the extraction procedure and not the data. One detail defeats the
obvious approach. Levels 1 to 4 each appear in two forms, 第N級 and 第N*級,
while levels 5 to 7 do not; filtering on the unstarred form --- which is what a
reader who has not noticed the asterisk will write --- silently discards 442 of
3{,}100 characters and 1{,}398 of 14{,}452 words, and produces totals that look
plausible. Counting both forms makes level 1 exactly 246 characters, the figure
NAER publishes, which is how we established that starred entries belong to
their numbered level. \textbf{What the asterisk denotes we do not determine.} It
does not track the word list's own 核心詞 core-word marking: both forms contain
core and non-core entries. Our extraction records the flag and leaves the
question open.

\section{Method}

Each Taiwan inventory is converted to simplified characters with OpenCC before
matching. Conversion is many-to-one, and 23 traditional entries collapse onto a
simplified form already present; we report this rather than absorb it. TBCL
writes alternates with a slash --- 爸爸/爸, 姊姊/姐姐/姊/姐 --- in 365 entries,
and we match each variant separately while counting the entry once.

\paragraph{Aligning the scales.} TBCL has seven levels. The mainland standards
have nine but grade 7, 8 and 9 as one undifferentiated band, so they have seven
distinct tiers. We align tier~7 to tier~7 as ``the top band'' and levels 1--6
one to one. \textbf{This is an assumption, not a finding.} No issuing body
publishes a TBCL-to-HSK correspondence, and \S\ref{sec:direction} shows how much
rests on it.

\paragraph{Control.} Before comparing across the strait we ran the same pipeline
over the two mainland documents, where the answer is published. It reproduces
\citet{serrano2026whichhsk} exactly: 9{,}674 shared words with 41.5\% at
different tiers, and 2{,}945 shared characters with 40.7\%. The cross-strait
figures below come from the same code path.

\section{Results}

\begin{table}[t]
\centering\small
\begin{tabular}{@{}lrrr@{}}
\toprule
\textbf{Comparison} & \textbf{Shared} & \textbf{Differ} & \textbf{Rate} \\
\midrule
\multicolumn{4}{@{}l}{\emph{Words}} \\
\quad PRC 2021 vs 2025   & 9{,}674 & 4{,}012 & 41.5\% \\
\quad TBCL vs PRC 2025   & 8{,}318 & 5{,}381 & \textbf{64.7\%} \\
\quad TBCL vs PRC 2021   & 7{,}885 & 5{,}607 & \textbf{71.1\%} \\
\addlinespace
\multicolumn{4}{@{}l}{\emph{Characters}} \\
\quad PRC 2021 vs 2025   & 2{,}945 & 1{,}200 & 40.7\% \\
\quad TBCL vs PRC 2025   & 2{,}785 & 1{,}538 & 55.2\% \\
\quad TBCL vs PRC 2021   & 2{,}737 & 1{,}601 & 58.5\% \\
\bottomrule
\end{tabular}
\caption{Shared items graded at different tiers.}
\label{tab:main}
\end{table}

The two ecosystems disagree with each other substantially more than the
mainland's own two documents do: 64.7\% of shared vocabulary against 41.5\%,
and 55.2\% of shared characters against 40.7\%. The gap is not explained by
inventory size, since TBCL and the 2025 syllabus share 8{,}318 words --- a
comparable base to the mainland pair's 9{,}674.

A practical reading: a text levelled under one ecosystem carries almost no
information about its level under the other. Two thirds of the vocabulary that
determines the answer is graded differently.

\section{The direction does not survive}
\label{sec:direction}

\begin{table}[t]
\centering\small
\begin{tabular}{@{}lrrr@{}}
\toprule
\textbf{Alignment} & \textbf{Differ} & \textbf{PRC} & \textbf{TBCL} \\
 & & \textbf{harder} & \textbf{harder} \\
\midrule
1:1 (as reported) & 64.7\% & 4{,}783 & 598 \\
TBCL $-1$ level   & 91.2\% & 7{,}497 & 85 \\
TBCL $-2$ levels  & 96.2\% & 7{,}982 & 19 \\
TBCL $+1$ level   & \textbf{43.9\%} & 1{,}234 & \textbf{2{,}421} \\
\bottomrule
\end{tabular}
\caption{Shifting one scale against the other. The rate never falls below the
mainland's internal 41.5\%; the direction reverses.}
\label{tab:sensitivity}
\end{table}

Under the 1:1 alignment the mainland syllabus appears to grade shared
vocabulary harder by 4{,}783 words to 598 --- a margin that invites a
conclusion about pedagogical conservatism. We do not draw it.

TBCL levels 1 and 2 sit below CEFR A1 in listening, reading and writing, while
the mainland scale describes no sub-A1 tier. If the scales are offset by a
level, a 1:1 alignment manufactures a direction. Table~\ref{tab:sensitivity}
shows that shifting TBCL by one level in the other direction reverses it
entirely, to 1{,}234 against 2{,}421.

What survives is the rate. Under every alignment tested it stays between 43.9\%
and 96.2\%, and its minimum still exceeds the 41.5\% disagreement inside the
mainland pair. \textbf{The finding is that the two ecosystems disagree; which
one is harder is not determinable without a correspondence that does not
exist.}

That absence is itself the result worth carrying. TOCFL score reports have
carried a TBCL column since August 2021, so Taiwan publishes an
exam-to-standard correspondence internally. Nothing comparable exists across
the strait, and the equivalency tables in circulation are reconstructions.

\section{Limitations}

The level alignment is the largest, and \S\ref{sec:direction} is the response to
it rather than a defence. Simplification is many-to-one and 23 entries collapse;
a comparison in the traditional direction would lose more, since the mainland
lists are published simplified. Matching is by orthographic form, so a word
shared in writing but differing in sense or register counts as shared --- the
figures are therefore a lower bound on disagreement, not an upper one. TOCFL is
not compared item by item here; we use it only for the count in
\S\ref{sec:counts}. TBCL's starred levels are counted at their numbered level on
the evidence of one published total, and a different reading of the asterisk
would change the per-level figures though not the totals. Finally, both Taiwan
inventories were extracted by us and validated against published totals, but not
against an official machine-readable release, because none exists.

\section*{Availability}

\texttt{scripts/tbcl\_extract.py} and \texttt{scripts/crossstrait.py} in the
\texttt{hsk30} repository reproduce every figure here from spreadsheets
downloaded free from the issuing bodies. Neither Taiwan inventory is
redistributed: NAER and SC-TOP assert rights over them, so the code ships and
the data does not. The mainland lists, the grading library and the corpus are
released as described in \citet{serrano2026whichhsk}.

\section*{Ethics and data statement}

No personal data and no human subjects. All four inventories are published
national reference documents, used here to compare those documents with each
other. Generative AI was used in writing the extraction and comparison code and
in drafting this paper; the problem framing, the design decisions and every
claim are the author's, who reviewed all output. That review is what identified
the alignment artifact reported in \S\ref{sec:direction}.

\bibliography{refs}

\end{document}
